\section{The Mandelbrot set and escape-time rendering}
\label{sec:mandelbrot-intro}

Escape-time rendering converts a question about an infinite sequence into a finite
pixel calculation. The distinction matters: observing escape can establish that a
parameter lies outside the set, while failing to observe escape within a budget
usually establishes only that more iterations may be needed.

The Mandelbrot set \(M\subset\mathbb{C}\) \cite{Mandelbrot1980Fractal} is defined as the set of complex
parameters \(c\) for which the orbit of
\begin{equation}
z_{n+1} = z_n^2 + c,\qquad z_0 = 0
\end{equation}
remains bounded. A common rendering method is \emph{escape-time}: iterate until
the orbit magnitude exceeds a bailout threshold, or until a maximum iteration
count is reached.

The standard mathematical escape guarantee uses \(|z_n|^2>4\), since
\(|z|>2\) implies divergence:
\begin{equation}
|z_n|^2 = \Re(z_n)^2 + \Im(z_n)^2 > 4 \quad\Rightarrow\quad \text{escaped.}
\end{equation}
For each pixel, the stored value is the smallest \(n\) (or a bounded proxy) at
which the selected rendering bailout is observed, or the maximum iteration limit
if no escape is observed within that budget. Kernels can use different bailout
radii and counting conventions, as described in their algorithm sections.

For $c=-1$, the orbit is $0,-1,0,-1,\ldots$. Returning to the same state proves
that this particular orbit repeats and is bounded. For $c=1$, the orbit begins
$0,1,2,5$ and escapes at the third update. A budget of two updates would not yet
observe that escape, illustrating why finite non-escape alone is inconclusive.
Equality at radius two is not itself an escape theorem; \cref{sec:real-form}
gives a bounded counterexample.


\subsection*{System architecture}
\label{sec:system-architecture}

At a high level, \FractalShark{}'s rendering pipeline transforms a view
specification (center, zoom, iteration limit) into a colored image through
a sequence of stages.  Most stages are implemented as CUDA kernels; the CPU
orchestrates data flow between them and handles high-precision coordinate
arithmetic.

\begin{lstlisting}[language={},basicstyle=\ttfamily\small,columns=flexible]
  View parameters (center, zoom, iterations)
          |
          v
  Coordinate mapping            [PointZoomBBConverter]
   (high-precision -> per-pixel deltas)
          |
          v
  Reference orbit computation   [RefOrbitCalc]
   (CPU: MPIR single/multi-thread,
    GPU: HpSharkFloat NTT pipeline)
          |
          v
  Linear approximation setup    [LAReference]
   (period detection, LA/BLA coefficient tables)
          |
          v
  Per-pixel perturbation (GPU)  [mandel_1xHDR_float_perturb_lav2]
   (LA stage traversal -> perturbation tail -> escape test)
          |
          v
  Antialiasing + coloring (GPU) [antialiasing_kernel]
   (supersample averaging, palette lookup)
          |
          v
  Display (OpenGL)              [RenderThreadPool GL consumer]
\end{lstlisting}

The first two stages---coordinate mapping and reference orbit
computation---dominate wall-clock time at extreme zoom depths, because
the required arithmetic precision grows with the logarithm of the zoom
factor.  Once the reference orbit is available, per-pixel perturbation
on the GPU is comparatively fast: each pixel evolves only a low-precision
delta relative to the shared reference
(\cref{sec:perturbation-concept}), with linear approximation
(\cref{sec:la-v2-perturb}) skipping the majority of iterations.

Two rendering paths exist (\cref{subsec:two-render-paths}): the
\emph{render pool path}, where worker threads each acquire a private
iteration buffer and produce frames for the OpenGL consumer; and the
\emph{direct path}, used by tests and save operations, which writes
results into the shared \code{m\_CurIters} buffer.  The render thread
pool (\cref{sec:render-thread-pool}) manages a command queue that
serializes UI-driven mutations with rendering work, discarding stale
renders when the user navigates faster than the GPU can complete them.

\subsection*{Document roadmap}

The remainder of this document is organized as follows:
\begin{itemize}
  \item \textbf{Direct rendering kernels and numeric types.}
        The GPU escape-time kernels and the fixed-point and floating-point
        numeric representations they operate on
        (\cref{sec:mandel-base-kernels,sec:hdr-float,sec:extended-precision-types}).
  \item \textbf{Perturbation theory and approximation-based acceleration.}
        Scaled perturbation, bilinear approximation (BLA), and the LA~v2
        pipeline that together enable deep-zoom rendering from a single
        reference orbit
        (\cref{sec:perturbation-concept,sec:bilinear-approx,sec:la-v2-perturb}).
  \item \textbf{Reference orbit computation.}
        CPU and GPU backends for computing the high-precision reference orbits
        required by perturbation theory (\cref{sec:ref-orbit-calc}).
  \item \textbf{GPU high-precision arithmetic.}
        NTT-based multiplication, parallel carry propagation, and the
        memory layout underlying GPU-side arbitrary-precision computation
        (\cref{subsec:ref-orbit-gpu,sec:hp-add,sec:ntt-multiply}).
  \item \textbf{Feature Finder.}
        Periodic-point detection via Newton's method, period scanning, and
        automatic deep-zoom target selection (\cref{sec:feature-finder}).
  \item \textbf{Rendering pipeline and memory infrastructure.}
        Antialiasing, coloring, file-backed growable vectors, and the overall
        render orchestration (\cref{sec:misc-details,sec:render-thread-pool,sec:disk_backed_growable_vectors}).
  \item \textbf{Appendices.}
        User documentation, a map of the source and build targets
        (\cref{sec:code-organization}), engineering reference material, and a
        development history recording implementation milestones and design
        decisions.
\end{itemize}

With the overview in place, Part~I turns to the rendering kernels themselves,
beginning with pixel-to-parameter mapping and the base escape-time iteration
that all subsequent optimizations build upon.
